\section{Identify the list codec before proving more list lemmas}
\label{sec:encoding}

\subsection{The remaining duplication}

The wire-format module \code{AnnotatedCodeArithmetic} defines
\code{listCode} on lists of natural numbers by
\[
  C([])=0,\qquad C(x::xs)=\operatorname{Nat.pair}(x,C(xs))+1.
\]
At the pinned Mathlib version, this is the same function as the installed
\code{Encodable} instance for \code{List Nat}. The equations arise in
different definitions, so the useful interface is a proved pointwise
identity, not a change of the repository's definition \cite{codearith,mathlist}.

The distinction is valuable because the existing code is not merely a
mathematical encoder. Supplied arithmetic programs, program-size statements,
and downstream reductions refer to its literal implementation. Replacing the
definition by a library alias would create an unnecessary definitional
compatibility obligation. Adding one theorem leaves those data untouched.

The actual bridge is short:
\begin{lean}
theorem listCode_eq_encode (xs : List Nat) :
    listCode xs = Encodable.encode xs := by
  induction xs with
  | nil => rfl
  | cons x xs ih => exact congrArg (fun n => Nat.pair x n + 1) ih
\end{lean}

The integration file keeps this as an explicitly invoked bridge. It does not
globally redirect every simplifier traversal of executable code into the
encoding typeclass.

\subsection{First consequence: reuse the standard length theorem}

\code{AnnotatedRename.length_le_tableCode}, in the module
\code{AnnotatedRenameCode}, currently proves its length
bound by a separate induction using a pairing inequality. A table is a list
of row codes, so the bridge reduces the entire argument to the library's
\code{Encodable.length_le_encode}:
\begin{lean}
  simpa only [tableCode, listCode_eq_encode, List.length_map] using
    (Encodable.length_le_encode (rows.map rowCode))
\end{lean}

This is stronger maintenance simplification than merely shortening the
induction. There is now one source for the general claim that a list's
length is bounded by this exact code. The row encoder requires no
injectivity hypothesis for this application: the theorem concerns the
length of the mapped list, and \code{List.length_map} supplies that equality
\cite{renamecode,mathlist}.

\subsection{Second consequence: derive typed parent checking from raw checking}

\code{AnnotatedParents} has both a typed-list theorem,
\code{check_lists}, and a raw-code theorem, \code{check_raw_iff}. The former
originally performs another induction over the expected list, with nested
fuel and parent-list cases. The latter already reconstructs the complete
parent spine from an arbitrary natural code. Its proof does not need
\code{check_lists} \cite{parents}.

Move \code{check_lists} below \code{check_raw_iff}, preserving its public
statement, and specialize the stronger theorem to \code{listCode parents}.
The only remaining issue is uniqueness of the decoded spine. The library
encoding is injective, and the bridge provides precisely that fact for
\code{listCode}:
\begin{lean}
  rw [check_raw_iff fuel history expected (listCode parents) hf]
  simp [listCode_eq_encode, Encodable.encode_injective.eq_iff]
\end{lean}

Schematically, the raw theorem has the form
\[
\mathrm{check}(c)=1
\quad\Longleftrightarrow\quad
\exists ps,\ c=C(ps)\;\land\;\mathrm{ParentsAgree}(ps).
\]
Setting \(c=C(parents)\) and applying injectivity removes the existential.
The premise \code{expected.length ≤ fuel} is retained exactly. No stronger
fuel assumption, parser termination argument, or exclusion of malformed
codes has been inserted.

\subsection{Why this survives the prior proposals}

Earlier reports propose lookup-library reuse, generic decoder graphs, and
using a decoder left inverse for code injectivity. Those are already
assumed available. The additional fact here is an exact identification with
an existing \emph{library instance}, which simultaneously supplies a length
bound and injectivity. The second change makes the typed result a corollary
of the independently proved raw result.
It does not introduce another scanner or decoder-correctness framework.

\subsection{Measured effect and integration order}

With all three modules included, the aggregate physical source length decreases
from 461 to 453 lines, counting the new import, explanatory comment, and
bridge theorem. The selected proof bodies go from 19 to 7 lines when the
bridge proof is included. All three modified modules were compiled with
their production imports.

There is one foundation-level tradeoff. The original table-length theorem
depends on \code{propext} and \code{Quot.sound}; the library replacement also
depends on \code{Classical.choice}, through the proof-bearing list encoding
instance and its library lemmas. This is a standard Lean axiom, not a new
assumption added to the project, but it is an observable change to
\code{#print axioms}. The parent-checker theorem already has all three
standard dependencies and does not increase its footprint. The table-length
change is therefore supplied as a separate optional patch. Omitting it
retains the original footprint and gives a four-line aggregate reduction
for the codec bridge plus parent specialization alone.

Apply \code{01-list-encoding-core.patch}, which includes the codec bridge
and parent theorem reordering. Apply \code{02-optional-table-length.patch}
only if the extra axiom dependency is acceptable. Existing public signatures
are unchanged.
The remaining maintenance obligation is explicit: if the repository's
encoding or Mathlib's chosen list encoding changes, the small bridge is
where that incompatibility will surface. It should not be hidden by an
overly broad simplification rule.
